\section{CONCLUSÃO}

Este trabalho confirmou a viabilidade técnica da predição neuromórfica de qualidade
da água no Lago do Serpa, combinando modelagem multimodal, análise espacial e
implantação embarcada. Em nível de desempenho, a SNN apresentou resultados
competitivos e estáveis, superando o MLP no benchmark multissemente e mantendo
proximidade de desempenho com baselines clássicos. No conjunto avaliado, Random
Forest e LSTM apresentaram menor erro médio absoluto global, enquanto a SNN
permaneceu como alternativa robusta para inferência neuromórfica em borda.
Além disso, o ajuste de pesos da função de perda mostrou melhora específica em
turbidez com trade-off moderado no desempenho macro.

No eixo de aplicação, o modelo foi exportado para TensorFlow Lite e integrado em
aplicativo móvel offline, com inferência local de pH, OD e turbidez, cálculo de
índice de poluição e visualização geoespacial com mapa offline (MBTiles). Esse
resultado avança o trabalho além de uma prova de conceito puramente algorítmica,
demonstrando um fluxo completo de dados, treinamento, inferência embarcada e apoio
à decisão em campo.
No desenho neuromórfico, a adoção de \textit{rate coding} também se mostrou coerente
com o cenário de borda, pois permite ajustar explicitamente o equilíbrio entre qualidade
da representação temporal e latência de inferência no celular por meio do parâmetro $T$. Como limitações,
o estudo ainda depende de base amostral local e reduzida, o que restringe a generalização para outros
corpos hídricos e regimes sazonais; além disso, o aumento de dados sintético não substitui coletas reais,
e diferenças entre modelos podem variar com novas campanhas de campo. A validação prática foi feita em
Android, Raspberry Pi e ESP32, mas a extrapolação para outros microcontroladores e diferentes perfis de
energia requer medições adicionais e campanhas de campo mais amplas. Como continuidade desta pesquisa,
recomenda-se ampliar campanhas de coleta, validar o modelo em múltiplos cenários ambientais e medir de
forma sistemática latência, consumo energético e robustez operacional em diferentes dispositivos de borda.
Também é relevante evoluir a calibração do índice composto de poluição com apoio de especialistas
ambientais e testar estratégias de atualização de modelo em campo para diferentes regimes hidrológicos.

Como contribuição adicional, a validação em ESP32 mostrou que a inferência multimodal é viável
em microcontroladores, com latência de ordem de milissegundos e erros absolutos compatíveis com
o cenário móvel. A comparação energética entre ESP32, Raspberry Pi e smartphones evidenciou
trade-offs claros entre throughput e energia por inferência, reforçando a importância de escolher
plataforma e regime de amostragem conforme o contexto operacional.
